\documentclass[../DoAn.tex]{subfiles}
\begin{document}

Phụ lục này đặc tả bổ sung một số use case đã được phân rã ở Chương~\ref{chapter:Related_works} nhưng không trình bày chi tiết trong phần nội dung chính. Các use case này thuộc nhóm chức năng quản trị nội dung, quản trị người dùng và vận hành hệ thống.

\section{Đặc tả use case ``Quản lý khóa học, chương và bài giảng''}

\begin{table}[H]
\centering
\small
\renewcommand{\arraystretch}{1.3}
\begin{tabular}{|>{\bfseries\raggedright\arraybackslash}p{2.6cm}|p{11.2cm}|}
\hline
Tên use case & Quản lý khóa học, chương và bài giảng \\ \hline
Tác nhân & Giảng viên \\ \hline
Mục đích & Cho phép giảng viên tổ chức nội dung đào tạo theo cấu trúc khóa học, chương và bài giảng; đồng thời quản lý trạng thái xuất bản của khóa học. \\
\hline
Tiền điều kiện & Giảng viên đã đăng nhập và có quyền quản lý trong tổ chức hoặc khóa học tương ứng. \\ \hline
Kịch bản chính & \begin{tabular}{@{}>{\raggedright\arraybackslash}p{0.9cm}>{\raggedright\arraybackslash}p{2.0cm}p{7.3cm}@{}}
\textbf{STT} & \textbf{Thực hiện bởi} & \textbf{Hành động} \\[2pt]
1 & Giảng viên & Tạo mới hoặc mở một khóa học trong tổ chức. \\[2pt]
2 & Giảng viên & Cập nhật thông tin khóa học như tiêu đề, mô tả, trạng thái và ảnh đại diện nếu có. \\[2pt]
3 & Giảng viên & Thực hiện quản lý chương (tạo mới, cập nhật tên, sắp xếp thứ tự hiển thị hoặc xóa chương). \\[2pt]
4 & Giảng viên & Thực hiện quản lý bài giảng trong từng chương (tạo mới, cập nhật thông tin bài giảng, sắp xếp thứ tự hiển thị hoặc xóa bài giảng). \\[2pt]
5 & Hệ thống & Lưu cấu trúc khóa học, chương và bài giảng; cập nhật danh sách hiển thị cho người dùng có quyền truy cập. \\[2pt]
\end{tabular} \\ \hline
Kịch bản khác (Luồng sự kiện thay thế) & \begin{tabular}{@{}>{\raggedright\arraybackslash}p{0.9cm}>{\raggedright\arraybackslash}p{2.0cm}p{7.3cm}@{}}
\textbf{STT} & \textbf{Thực hiện bởi} & \textbf{Hành động} \\[2pt]
5a & Giảng viên & Thay đổi trạng thái xuất bản của khóa học; hệ thống cập nhật lại danh sách khóa học hiển thị phía người học. \\[2pt]
\end{tabular} \\ \hline
Hậu điều kiện & Cấu trúc khóa học được lưu và sẵn sàng cho bước tải video bài giảng hoặc xuất bản tới người học. \\ \hline
\end{tabular}
\caption{Đặc tả use case Quản lý khóa học, chương và bài giảng}
\label{table:pl_uc_quanlykhoahoc}
\end{table}

\section{Đặc tả use case ``Tải lên video bài giảng''}

\begin{table}[H]
\centering
\small
\renewcommand{\arraystretch}{1.3}
\begin{tabular}{|>{\bfseries\raggedright\arraybackslash}p{2.6cm}|p{11.2cm}|}
\hline
Tên use case & Tải lên video bài giảng \\ \hline
Tác nhân & Giảng viên \\ \hline
Mục đích & Cho phép giảng viên tải video bài giảng lên hệ thống từ thiết bị cá nhân để chuẩn bị cho quá trình học tập và phân tích sinh tương tác. \\ \hline
Tiền điều kiện & Giảng viên đã đăng nhập và bài giảng tương ứng đã được khởi tạo. \\ \hline
Kịch bản chính & \begin{tabular}{@{}>{\raggedright\arraybackslash}p{0.9cm}>{\raggedright\arraybackslash}p{2.0cm}p{7.3cm}@{}}
\textbf{STT} & \textbf{Thực hiện bởi} & \textbf{Hành động} \\[2pt]
1 & Giảng viên & Mở giao diện quản lý bài giảng, chọn chức năng tải lên video cho bài giảng tương ứng. \\[2pt]
2 & Hệ thống & Hiển thị khu vực chọn tệp và các thông số giới hạn (định dạng mp4, dung lượng tối đa). \\[2pt]
3 & Giảng viên & Chọn tệp video từ thiết bị cá nhân và nhấn nút tải lên. \\[2pt]
4 & Hệ thống & Tiến hành truyền tải video lên kho lưu trữ đối tượng (S3) và hiển thị thanh tiến độ tải lên. \\[2pt]
5 & Hệ thống & Sau khi tải lên thành công, lưu liên kết video vào cơ sở dữ liệu bài giảng và thông báo hoàn tất cho giảng viên. \\[2pt]
\end{tabular} \\ \hline
Kịch bản khác (Luồng sự kiện thay thế) & \begin{tabular}{@{}>{\raggedright\arraybackslash}p{0.9cm}>{\raggedright\arraybackslash}p{2.0cm}p{7.3cm}@{}}
\textbf{STT} & \textbf{Thực hiện bởi} & \textbf{Hành động} \\[2pt]
3a & Hệ thống & Nếu tệp tải lên không phải video hợp lệ hoặc vượt giới hạn dung lượng, hệ thống từ chối tải lên và hiển thị thông báo lỗi. \\[2pt]
4a & Hệ thống & Nếu kết nối mạng bị gián đoạn giữa chừng, dừng tiến trình và hiển thị thông báo lỗi tải lên để giảng viên thực hiện lại. \\[2pt]
\end{tabular} \\ \hline
Hậu điều kiện & Tệp video bài giảng được lưu trữ thành công trên máy chủ và sẵn sàng để kích hoạt phân tích. \\ \hline
\end{tabular}
\caption{Đặc tả use case Tải lên video bài giảng}
\label{table:pl_uc_tailenvideo}
\end{table}

\section{Đặc tả use case ``Biên tập bản chép lời''}

\begin{table}[H]
\centering
\small
\renewcommand{\arraystretch}{1.3}
\begin{tabular}{|>{\bfseries\raggedright\arraybackslash}p{2.6cm}|p{11.2cm}|}
\hline
Tên use case & Biên tập bản chép lời \\ \hline
Tác nhân & Giảng viên \\ \hline
Mục đích & Cho phép giảng viên hiệu chỉnh bản chép lời và các đoạn nội dung được sinh tự động trước khi dùng làm cơ sở tạo hoạt động tương tác. \\ \hline
Tiền điều kiện & Bài giảng đã được phân tích hoặc đã có bản chép lời/phân đoạn nội dung. \\ \hline
Kịch bản chính & \begin{tabular}{@{}>{\raggedright\arraybackslash}p{0.9cm}>{\raggedright\arraybackslash}p{2.0cm}p{7.3cm}@{}}
\textbf{STT} & \textbf{Thực hiện bởi} & \textbf{Hành động} \\[2pt]
1 & Giảng viên & Mở màn hình phân tích của bài giảng. \\[2pt]
2 & Hệ thống & Hiển thị các đoạn nội dung theo thứ tự thời gian, kèm mốc bắt đầu, mốc kết thúc và tóm tắt. \\[2pt]
3 & Giảng viên & Sửa nội dung lời thoại để khắc phục lỗi nhận dạng. \\[2pt]
4 & Giảng viên & Gộp hoặc tách phân đoạn nội dung. \\[2pt]
5 & Giảng viên & Điều chỉnh ranh giới thời gian bắt đầu và kết thúc của phân đoạn. \\[2pt]
6 & Hệ thống & Kiểm tra tính hợp lệ của dữ liệu và lưu thay đổi. \\[2pt]
\end{tabular} \\ \hline
Kịch bản khác (Luồng sự kiện thay thế) & \begin{tabular}{@{}>{\raggedright\arraybackslash}p{0.9cm}>{\raggedright\arraybackslash}p{2.0cm}p{7.3cm}@{}}
\textbf{STT} & \textbf{Thực hiện bởi} & \textbf{Hành động} \\[2pt]
3a & Giảng viên & Thêm một phân đoạn thủ công khi bản chép lời thiếu đoạn quan trọng. \\[2pt]
3b & Giảng viên & Xóa phân đoạn không phù hợp; hệ thống cập nhật lại thứ tự. \\[2pt]
6a & Hệ thống & Nếu mốc thời gian không hợp lệ hoặc dữ liệu rỗng, hệ thống không lưu và yêu cầu chỉnh sửa lại. \\[2pt]
\end{tabular} \\ \hline
Hậu điều kiện & Bản chép lời và các đoạn nội dung được cập nhật, bảo đảm phù hợp hơn cho bước sinh hoạt động tương tác. \\ \hline
\end{tabular}
\caption{Đặc tả use case Biên tập bản chép lời}
\label{table:pl_uc_bientapbanchep}
\end{table}

\section{Đặc tả use case ``Quản lý thành viên và duyệt yêu cầu tham gia''}

\begin{table}[H]
\centering
\small
\renewcommand{\arraystretch}{1.3}
\begin{tabular}{|>{\bfseries\raggedright\arraybackslash}p{2.6cm}|p{11.2cm}|}
\hline
Tên use case & Quản lý thành viên và duyệt yêu cầu tham gia \\ \hline
Tác nhân & Giảng viên \\ \hline
Mục đích & Cho phép giảng viên kiểm soát danh sách học viên trong khóa học và xử lý các yêu cầu tham gia khóa học. \\ \hline
Tiền điều kiện & Khóa học đã tồn tại; giảng viên có quyền quản lý khóa học. \\ \hline
Kịch bản chính & \begin{tabular}{@{}>{\raggedright\arraybackslash}p{0.9cm}>{\raggedright\arraybackslash}p{2.0cm}p{7.3cm}@{}}
\textbf{STT} & \textbf{Thực hiện bởi} & \textbf{Hành động} \\[2pt]
1 & Giảng viên & Mở danh sách yêu cầu tham gia khóa học. \\[2pt]
2 & Hệ thống & Hiển thị các yêu cầu đang chờ duyệt cùng thông tin người học. \\[2pt]
3 & Giảng viên & Chấp nhận một yêu cầu tham gia. \\[2pt]
4 & Hệ thống & Cập nhật trạng thái yêu cầu thành đã duyệt và thêm người học vào danh sách thành viên khóa học. \\[2pt]
5 & Giảng viên & Xem hoặc cập nhật danh sách thành viên hiện tại của khóa học. \\[2pt]
\end{tabular} \\ \hline
Kịch bản khác (Luồng sự kiện thay thế) & \begin{tabular}{@{}>{\raggedright\arraybackslash}p{0.9cm}>{\raggedright\arraybackslash}p{2.0cm}p{7.3cm}@{}}
\textbf{STT} & \textbf{Thực hiện bởi} & \textbf{Hành động} \\[2pt]
3a & Giảng viên & Từ chối yêu cầu; hệ thống cập nhật trạng thái bị từ chối và không thêm người học vào khóa học. \\[2pt]
5a & Giảng viên & Thêm trực tiếp một người học là thành viên tổ chức vào khóa học. \\[2pt]
5b & Giảng viên & Xóa một thành viên khỏi khóa học khi người đó không còn tham gia lớp. \\[2pt]
\end{tabular} \\ \hline
Hậu điều kiện & Danh sách thành viên khóa học và trạng thái yêu cầu tham gia được cập nhật nhất quán. \\ \hline
\end{tabular}
\caption{Đặc tả use case Quản lý thành viên và duyệt yêu cầu tham gia}
\label{table:pl_uc_thanhvien}
\end{table}

\section{Đặc tả use case ``Xem kết quả và phản hồi''}

\begin{table}[H]
\centering
\small
\renewcommand{\arraystretch}{1.3}
\begin{tabular}{|>{\bfseries\raggedright\arraybackslash}p{2.6cm}|p{11.2cm}|}
\hline
Tên use case & Xem kết quả và phản hồi \\ \hline
Tác nhân & Người học \\ \hline
Mục đích & Cho phép người học xem lại điểm, câu trả lời, phản hồi chi tiết và theo dõi tiến độ hoàn thành của bản thân trong khóa học. \\ \hline
Tiền điều kiện & Người học đã làm ít nhất một hoạt động tương tác hoặc đã nộp bài cho bài giảng. \\ \hline
Kịch bản chính & \begin{tabular}{@{}>{\raggedright\arraybackslash}p{0.9cm}>{\raggedright\arraybackslash}p{2.0cm}p{7.3cm}@{}}
\textbf{STT} & \textbf{Thực hiện bởi} & \textbf{Hành động} \\[2pt]
1 & Người học & Mở bài giảng hoặc trang kết quả học tập cá nhân. \\[2pt]
2 & Hệ thống & Tải các lần làm bài, điểm số, tiến độ xem video và phản hồi đã lưu. \\[2pt]
3 & Người học & Chọn xem bảng theo dõi tiến độ học tập tổng quát của toàn khóa học. \\[2pt]
4 & Người học & Chọn một lần làm bài cụ thể của bài giảng để xem chi tiết. \\[2pt]
5 & Hệ thống & Hiển thị từng câu trả lời, điểm đạt được và nhận xét tương ứng của lần làm bài đó. \\[2pt]
6 & Người học & Dựa trên phản hồi để xem lại đoạn video bài giảng hoặc thực hiện lượt làm bài tiếp theo nếu còn hạn. \\[2pt]
\end{tabular} \\ \hline
Kịch bản khác (Luồng sự kiện thay thế) & \begin{tabular}{@{}>{\raggedright\arraybackslash}p{0.9cm}>{\raggedright\arraybackslash}p{2.0cm}p{7.3cm}@{}}
\textbf{STT} & \textbf{Thực hiện bởi} & \textbf{Hành động} \\[2pt]
2a & Hệ thống & Nếu chưa có dữ liệu làm bài, hệ thống hiển thị giao diện trống và hướng dẫn người học bắt đầu bài học. \\[2pt]
5a & Hệ thống & Nếu cấu hình phản hồi của bài giảng chưa cho phép xem đáp án, hệ thống chỉ hiển thị điểm và nhận xét phù hợp. \\[2pt]
\end{tabular} \\ \hline
Hậu điều kiện & Người học nắm được kết quả học tập cá nhân và tiến độ hoàn thành khóa học để tiếp tục cải thiện. \\ \hline
\end{tabular}
\caption{Đặc tả use case Xem kết quả và phản hồi}
\label{table:pl_uc_xemketqua_app}
\end{table}

\section{Đặc tả use case ``Quản lý người dùng và phân quyền''}

\begin{table}[H]
\centering
\small
\renewcommand{\arraystretch}{1.3}
\begin{tabular}{|>{\bfseries\raggedright\arraybackslash}p{2.6cm}|p{11.2cm}|}
\hline
Tên use case & Quản lý người dùng và phân quyền \\ \hline
Tác nhân & Quản trị viên \\ \hline
Mục đích & Cho phép quản trị viên quản lý tài khoản người dùng, trạng thái tài khoản và vai trò hệ thống. \\ \hline
Tiền điều kiện & Quản trị viên đã đăng nhập bằng tài khoản có quyền quản trị toàn hệ thống. \\ \hline
Kịch bản chính & \begin{tabular}{@{}>{\raggedright\arraybackslash}p{0.9cm}>{\raggedright\arraybackslash}p{2.0cm}p{7.3cm}@{}}
\textbf{STT} & \textbf{Thực hiện bởi} & \textbf{Hành động} \\[2pt]
1 & Quản trị viên & Mở màn hình danh sách người dùng. \\[2pt]
2 & Hệ thống & Hiển thị danh sách người dùng và hỗ trợ tìm kiếm theo định danh, email hoặc tên phù hợp. \\[2pt]
3 & Quản trị viên & Tạo mới hoặc cập nhật thông tin tài khoản người dùng. \\[2pt]
4 & Quản trị viên & Gán vai trò hệ thống hoặc thay đổi trạng thái hoạt động của tài khoản. \\[2pt]
5 & Hệ thống & Lưu thay đổi và áp dụng quyền mới cho các lần truy cập sau. \\[2pt]
\end{tabular} \\ \hline
Kịch bản khác (Luồng sự kiện thay thế) & \begin{tabular}{@{}>{\raggedright\arraybackslash}p{0.9cm}>{\raggedright\arraybackslash}p{2.0cm}p{7.3cm}@{}}
\textbf{STT} & \textbf{Thực hiện bởi} & \textbf{Hành động} \\[2pt]
3a & Hệ thống & Nếu email đã tồn tại hoặc dữ liệu không hợp lệ, hệ thống từ chối tạo/cập nhật và hiển thị lỗi. \\[2pt]
4a & Quản trị viên & Khóa tài khoản; hệ thống không cho phép tài khoản đó đăng nhập hoặc truy cập chức năng yêu cầu quyền. \\[2pt]
\end{tabular} \\ \hline
Hậu điều kiện & Thông tin tài khoản, vai trò và trạng thái người dùng được cập nhật trong hệ thống. \\ \hline
\end{tabular}
\caption{Đặc tả use case Quản lý người dùng và phân quyền}
\label{table:pl_uc_nguoidung}
\end{table}

\section{Đặc tả use case ``Quản lý tổ chức và thành viên''}

\begin{table}[H]
\centering
\small
\renewcommand{\arraystretch}{1.3}
\begin{tabular}{|>{\bfseries\raggedright\arraybackslash}p{2.6cm}|p{11.2cm}|}
\hline
Tên use case & Quản lý tổ chức và thành viên \\ \hline
Tác nhân & Quản trị viên \\ \hline
Mục đích & Cho phép quản trị viên tạo, sửa, xóa các tổ chức đào tạo và quản trị danh sách cũng như vai trò của các thành viên trong tổ chức đó. \\ \hline
Tiền điều kiện & Quản trị viên đã đăng nhập hệ thống. \\ \hline
Kịch bản chính & \begin{tabular}{@{}>{\raggedright\arraybackslash}p{0.9cm}>{\raggedright\arraybackslash}p{2.0cm}p{7.3cm}@{}}
\textbf{STT} & \textbf{Thực hiện bởi} & \textbf{Hành động} \\[2pt]
1 & Quản trị viên & Mở giao diện quản lý tổ chức đào tạo. \\[2pt]
2 & Quản trị viên & Thực hiện tạo mới hoặc điều chỉnh thông tin chung của tổ chức (tên, slug định danh). \\[2pt]
3 & Quản trị viên & Chọn xem danh sách thành viên thuộc tổ chức tương ứng. \\[2pt]
4 & Hệ thống & Tải và hiển thị danh sách thành viên kèm vai trò hiện tại của từng người. \\[2pt]
5 & Quản trị viên & Thực hiện thêm tài khoản mới vào tổ chức hoặc điều chỉnh vai trò thành viên (chủ sở hữu, quản trị, giảng viên, học viên). \\[2pt]
6 & Hệ thống & Lưu trữ cấu trúc tổ chức và phân quyền thành viên mới cập nhật. \\[2pt]
\end{tabular} \\ \hline
Kịch bản khác (Luồng sự kiện thay thế) & \begin{tabular}{@{}>{\raggedright\arraybackslash}p{0.9cm}>{\raggedright\arraybackslash}p{2.0cm}p{7.3cm}@{}}
\textbf{STT} & \textbf{Thực hiện bởi} & \textbf{Hành động} \\[2pt]
2a & Hệ thống & Nếu đường dẫn định danh (slug) của tổ chức bị trùng lặp, hệ thống từ chối lưu và yêu cầu chỉnh sửa lại. \\[2pt]
5a & Quản trị viên & Xóa thành viên ra khỏi tổ chức; hệ thống thu hồi toàn bộ quyền liên quan đến tổ chức đó. \\[2pt]
\end{tabular} \\ \hline
Hậu điều kiện & Thông tin cơ cấu tổ chức và phân quyền thành viên được cập nhật chính xác trên toàn hệ thống. \\ \hline
\end{tabular}
\caption{Đặc tả use case Quản lý tổ chức và thành viên}
\label{table:pl_uc_tochuc}
\end{table}

\section{Đặc tả use case ``Giám sát tác vụ phân tích''}

\begin{table}[H]
\centering
\small
\renewcommand{\arraystretch}{1.3}
\begin{tabular}{|>{\bfseries\raggedright\arraybackslash}p{2.6cm}|p{11.2cm}|}
\hline
Tên use case & Giám sát tác vụ phân tích \\ \hline
Tác nhân & Quản trị viên \\ \hline
Mục đích & Cho phép quản trị viên theo dõi các tác vụ phân tích video đang chạy, đã hoàn thành hoặc thất bại trên toàn hệ thống. \\ \hline
Tiền điều kiện & Hệ thống đã phát sinh ít nhất một tác vụ phân tích hoặc sinh nội dung tương tác. \\ \hline
Kịch bản chính & \begin{tabular}{@{}>{\raggedright\arraybackslash}p{0.9cm}>{\raggedright\arraybackslash}p{2.0cm}p{7.3cm}@{}}
\textbf{STT} & \textbf{Thực hiện bởi} & \textbf{Hành động} \\[2pt]
1 & Quản trị viên & Mở màn hình giám sát tác vụ. \\[2pt]
2 & Hệ thống & Hiển thị danh sách tác vụ kèm loại tác vụ, trạng thái, tiến độ, thời điểm tạo và người khởi tạo. \\[2pt]
3 & Quản trị viên & Mở chi tiết một tác vụ. \\[2pt]
4 & Hệ thống & Hiển thị thông tin bài giảng liên quan, bước xử lý hiện tại, thông báo lỗi nếu có và dữ liệu tiến độ. \\[2pt]
5 & Quản trị viên & Theo dõi trạng thái để phát hiện tác vụ lỗi hoặc tồn đọng bất thường. \\[2pt]
\end{tabular} \\ \hline
Kịch bản khác (Luồng sự kiện thay thế) & \begin{tabular}{@{}>{\raggedright\arraybackslash}p{0.9cm}>{\raggedright\arraybackslash}p{2.0cm}p{7.3cm}@{}}
\textbf{STT} & \textbf{Thực hiện bởi} & \textbf{Hành động} \\[2pt]
5a & Quản trị viên & Hủy một tác vụ đang chờ xử lý hoặc đang chạy; hệ thống cập nhật trạng thái tác vụ thành đã hủy. \\[2pt]
5b & Hệ thống & Nếu tác vụ thất bại, hệ thống lưu thông điệp lỗi để quản trị viên hoặc giảng viên có cơ sở xử lý lại. \\[2pt]
\end{tabular} \\ \hline
Hậu điều kiện & Trạng thái vận hành của các tác vụ phân tích được theo dõi tập trung, hỗ trợ phát hiện và xử lý lỗi trong quá trình vận hành. \\ \hline
\end{tabular}
\caption{Đặc tả use case Giám sát tác vụ phân tích}
\label{table:pl_uc_giamsat}
\end{table}

\end{document}
